
\documentclass[UTF8,11pt,a4paper]{ctexart}
\usepackage[a4paper,left=23mm,right=23mm,top=24mm,bottom=24mm]{geometry}
\usepackage{fontspec}
\setmainfont{DejaVu Sans}
\setCJKmainfont{Noto Serif CJK SC}
\usepackage{amsmath,amssymb}
\usepackage{xcolor}
\usepackage[colorlinks=true,linkcolor=blue!45!black,unicode=true]{hyperref}
\usepackage{enumitem}
\usepackage{needspace}
\usepackage{fancyhdr}
\setlength{\parindent}{2em}
\setlength{\parskip}{0.25em}
\setlist[enumerate]{leftmargin=2.6em,itemsep=0.4em,topsep=0.5em}
\pagestyle{fancy}
\fancyhf{}
\fancyhead[L]{\small 三千弱水 · 大学化学高阶综合题库}
\fancyhead[R]{\small 待专家审校稿}
\fancyfoot[C]{\thepage}
\setlength{\headheight}{15pt}
\sloppy

\begin{document}
\begin{center}
{\LARGE\bfseries 大学化学高阶综合题 · 第10辑 (16题)}\\[0.8em]
{\large 题目、分步解析与数值核验}\\[0.5em]
{\small \textcolor{red!65!black}{审校稿：未经化学专家核准；禁止直接视为正式试卷或出版物}}
\end{center}
\vspace{1em}
\noindent 本册从仓库结构化 JSONL 自动生成。增补和修改应只编辑源题目数据，并重新执行统一编译。
\par 机械方程残差检查只能说明指定模型的内部数值一致，不保证模型现实可行性。第六辑另须完成跨辑模型去重审查。
\tableofcontents
\clearpage
\section*{第10辑 · 题目及解答}
\addcontentsline{toc}{section}{第10辑}
\subsection*{无机化学}
\addcontentsline{toc}{subsection}{第10辑 / 无机化学}
\Needspace{7\baselineskip}
\subsubsection*{UC-ADV10-I101 · 反铁磁二核交换耦合的自旋统计与磁响应}
\noindent\textbf{题目}\par
假想二核过渡金属配合物中各中心自旋S1=S2=1/2，采用各向同性Heisenberg交换H=J S1·S2，J>0代表反铁磁。低能单态S=0简并度1，高能三重态S=1简并度3，能隙J=2.400kJ/mol。温度280.0K且达到热平衡，忽略激发轨道和强场自旋混合。求三重态和单态布居、磁化率相对两个不耦合S=1/2中心Curie磁化率的比值、平均激发能与摩尔热容，并验证高温磁化率的正确极限。\par
\medskip\noindent\textbf{分步解答}
\begin{enumerate}
\item 取单态能量为0，三重态能量为J；配分函数Z=1+3exp(−J/RT)。
\item 设q=exp(−J/RT)=0.3566836943，Z=2.070051083。
\item 三重态pT=3q/Z=0.5169201339，单态pS=1/Z=0.4830798661。
\item 小场磁化率按照S(S+1)加权，与两个自由S=1/2自旋比较得χ/χfree=4q/(1+3q)=0.6892268451。
\item 单态为参考零能时U=JpT=1240.608321J/mol。
\item 通过对U求T导数得到C=(J²/(RT²))pT(1−pT)=2.206554309J/(mol K)。
\item 高温q→1给χ/χfree→1，低温q→0给单态占据趋1；热容在低高温均趋零。
\item 数值回代核验残差：0.00e+0、0.00e+0、0.00e+0、0.00e+0。
\end{enumerate}
\noindent\textbf{答案摘要}\par
答案1 = 0.516920134 fraction；答案2 = 0.483079866 fraction；答案3 = 0.689226845 relative；答案4 = 1240.60832 J/mol；答案5 = 2.20655431 J/(mol K)\par
\noindent\textbf{模型前提：}
仅按题设理想化模型处理；实际化学性、跨辑核心模型独立性、版权和专家审校仍待核准。\par
\noindent\textcolor{red!65!black}{[跨辑去重未完成 · 草稿]} \par
\medskip\hrule\medskip
\Needspace{7\baselineskip}
\subsubsection*{UC-ADV10-I102 · 质子耦合电子转移方形热力学循环与pH形式电势}
\noindent\textbf{题目}\par
某假想无机金属配合物发生一电子氧化还原，氧化态OxH与还原态RedH可分别失去一个质子成为Ox和Red。已知pKa(OxH)=5.000、pKa(RedH)=8.000，在298.15K时质子化态一电子标准还原电位OxH+e−→RedH为E°H=0.6200V。忽略进一步质子化及其它配位变化，假设四种物种理想且氧化态与还原态总活度相等。求完全去质子化的Ox/Red一电子标准还原电势、pH=6.500时的形式电势、两态在该pH下的质子化分数，以及pH=2.000的形式电势，并说明高低pH极限。\par
\medskip\noindent\textbf{分步解答}
\begin{enumerate}
\item 写两个酸解离平衡：OxH⇌Ox+H+、RedH⇌Red+H+，各平衡常数分别10\textasciicircum{}−5与10\textasciicircum{}−8。
\item 由热力学方形循环，去质子化标准电位E°deprot=E°H+(RT/F)ln(Ka,Red/Ka,Ox)=0.442521951V；此处还原态酸性更弱，使去质子化电位变低。
\item pH=6.500下氧化态质子化份额αOxH=1/(1+10\textasciicircum{}(pH−pKa,Ox))=0.03065343。
\item 对应还原态质子化份额αRedH=1/(1+10\textasciicircum{}(pH−pKa,Red))=0.96934657。
\item 将Nernst式用于各自总活度相等的条件，Eformal=EH+(RT/F)ln(αOxH/αRedH)=0.531260975V。
\item pH=2.000时形式电位为0.619974346V，接近全质子化电位0.6200V；pH远大于8时趋于0.442521951V。
\item 平衡常数与形式电势在任意pH遵循同一热力学闭合循环；若有氧化态聚集或电子数不同，该四角模型失效。
\item 独立代回原始配分、强度或热力学公式所得主要残差：0.000e+0、-2.776e-17、0.000e+0；若超出题设近似则需重新建模。
\end{enumerate}
\noindent\textbf{答案摘要}\par
答案1 = 0.442521951 V；答案2 = 0.531260975 V；答案3 = 0.03065343 fraction；答案4 = 0.96934657 fraction；答案5 = 0.619974346 V\par
\noindent\textbf{模型前提：}
标准态、理想活度和模型参数仅按本题设定成立，未经专家科学性及跨辑原创性审定。\par
\noindent\textcolor{red!65!black}{[跨辑去重未完成 · 草稿]} \par
\medskip\hrule\medskip
\Needspace{7\baselineskip}
\subsubsection*{UC-ADV10-I103 · Born–Landé离子晶体静电晶格能与NaCl型密度校核}
\noindent\textbf{题目}\par
假想NaCl型立方一价离子晶体，Madelung常数A\_M=1.74756、最近邻正负离子距离r0=0.2820nm、Born排斥指数n=9.000。用ε0=8.8541878128×10\textasciicircum{}-12F/m、e=1.602176634×10\textasciicircum{}-19C、NA=6.02214076×10\textasciicircum{}23mol\textasciicircum{}-1计算每mol化学式单元晶格生成能U=−NA A\_M e²/(4πε0r0)(1−1/n)。另给假想摩尔质量58.44g/mol、常规NaCl型晶胞含4个化学式单元，几何晶格常数a=2r0。求不含短程排斥修正的库仑能量大小、晶格生成能、反向解离能、理论晶体密度和排斥校正所占库仑项百分比，并核对负号代表晶格形成放热。\par
\medskip\noindent\textbf{分步解答}
\begin{enumerate}
\item 不含Born短程排斥修正的库仑吸引大小Uc=NA A\_M e²/(4πε0r0)=860.985987kJ/mol。
\item Born–Landé晶格形成能U=−Uc(1−1/n)=-765.320877kJ/mol。
\item 气态自由离子分离固体的反向晶格解离能=−U=765.320877kJ/mol，正号代表需供能。
\item 排斥能修正相对Uc的比例为1/n=0.111111111（11.1111111\%），不能重复把晶格生成能换号。
\item NaCl常规晶胞边长a=2r0=0.564nm，含4个化学式单元。
\item 使用式量0.05844kg/mol，密度ρ=4MW/[NA(2r0)\textasciicircum{}3]=2163.62497kg/m³。
\item 真实晶格能还包含极化、色散及离子非球形效应；Born–Landé表达式只是所指定的离子晶体近似。
\item 独立代回原始配分、强度或热力学公式所得主要残差：0.000e+0、0.000e+0、-1.110e-16；若超出题设近似则需重新建模。
\end{enumerate}
\noindent\textbf{答案摘要}\par
答案1 = 860.985987 kJ/mol；答案2 = -765.320877 kJ/mol；答案3 = 765.320877 kJ/mol；答案4 = 2163.62497 kg/m³；答案5 = 0.111111111 fraction\par
\noindent\textbf{模型前提：}
标准态、理想活度和模型参数仅按本题设定成立，未经专家科学性及跨辑原创性审定。\par
\noindent\textcolor{red!65!black}{[跨辑去重未完成 · 草稿]} \par
\medskip\hrule\medskip
\Needspace{7\baselineskip}
\subsubsection*{UC-ADV10-I104 · 尖晶石反位阳离子分布的构型熵与平衡反位度}
\noindent\textbf{题目}\par
考虑假想正常尖晶石AB₂O₄，其中A为二价、B为三价，氧为二价；每化学式单位有1个四面体阳离子位与2个八面体阳离子位。设反位度x使四面体位组成为A\_(1−x)B\_x、八面体位组成为[A\_x B\_(2−x)]，0≤x≤1。每摩尔化学式单位的反位能代价为Δx，Δ=12.00 kJ/mol；T=1000.0 K，忽略相互作用并将两类位点视作各自理想混合。求平衡x、八面体位A和B的位点分数、每摩尔化学式单位构型熵S\_conf/R、相对x=0有序态的平衡自由能G−G₀，并说明Δ=0与低温极限。\par
\medskip\noindent\textbf{分步解答}
\begin{enumerate}
\item 按每个化学式单位计，四面体A、B分别为1−x与x，八面体A、B分别为x与2−x；总A=1、总B=2，离子电荷2+6−8=0。
\item 两类子晶格各自理想混合，构型熵除以R为−[(1−x)ln(1−x)+x ln x+x ln(x/2)+(2−x)ln((2−x)/2)]，其中0 ln 0取连续极限0。
\item 以x=0有序态为零，G(x)−G₀=Δx−T S\_conf(x)；驻点条件dG/dx=Δ+RT ln[x²/((1−x)(2−x))]=0。
\item 设K=exp(−Δ/RT)=0.2361546810，平衡方程x²=K(1−x)(2−x)，即(1−K)x²+3Kx−2K=0。
\item 选择0≤x≤1的根x=0.4491556286；另一个根为负，不能代表位点占据。二阶导数RT[2/x+1/(1−x)+1/(2−x)]>0，故驻点是唯一稳定极小值。
\item 八面体位A分数x/2=0.2245778143，B分数1−x/2=0.7754221857；四面体A、B分数之和也为1。
\item 代回理想混合熵得S\_conf/R=1.7532503045，代回G(x)−G₀=Δx−RT(S\_conf/R)=−9187.466573 J/mol；自由能负值说明在该温度熵增足以抵消反位能。
\item 若Δ=0，随机分布给x=2/3；若Δ>0且T→0，则K→0、x→0。模型未考虑真实尖晶石晶场能的非线性、振动熵及位点相关性。
\end{enumerate}
\noindent\textbf{答案摘要}\par
答案1 = 0.449155629 fraction；答案2 = 0.224577814 fraction；答案3 = 0.775422186 fraction；答案4 = 1.7532503 dimensionless；答案5 = -9187.46657 J/mol\par
\noindent\textbf{模型前提：}
AB₂O₄仅为理想化二价/三价尖晶石；反位能线性、理想子晶格混合、定温平衡；跨辑模型和真实材料适用性仍待人工复核。\par
\noindent\textcolor{red!65!black}{[跨辑去重未完成 · 草稿]} \par
\medskip\hrule\medskip
\subsection*{分析化学}
\addcontentsline{toc}{subsection}{第10辑 / 分析化学}
\Needspace{7\baselineskip}
\subsubsection*{UC-ADV10-A101 · DOSY脉冲梯度NMR的Stejskal–Tanner衰减与自扩散反演}
\noindent\textbf{题目}\par
一份假想溶液经脉冲梯度NMR的DOSY法分析，质子旋磁比γ=2.67522×10\textasciicircum{}8rad/(s T)，梯度幅度g=0.2500T/m，脉冲宽度δ=2.500ms，扩散等待时间Δ=45.00ms，样品自扩散系数D=6.000×10\textasciicircum{}-10m²/s。给定Stejskal–Tanner式I/I0=exp(−bD)，b=γ²g²δ²(Δ−δ/3)。求梯度系数b、回波归一化强度、从强度逆推的D，以及将梯度g加倍后的回波强度，指出为何指数不能只加倍。\par
\medskip\noindent\textbf{分步解答}
\begin{enumerate}
\item 将脉冲时间统一为秒，δ=0.002500s、Δ=0.04500s。
\item 梯度因子b=γ²g²δ²(Δ−δ/3)=1234734728s/m²。
\item 衰减指数bD=0.740840837，故归一化强度s=exp(−bD)=0.4767129091。
\item 由观测强度反演D=−ln(s)/b=6e-10m²/s。
\item 若梯度幅度加倍，b变为4b，而不是2b；对应s2=exp(−4bD)=0.05164492488。
\item b乘D无量纲；γ应使用弧度/s/T，否则需正确处理2π因子。
\item 该模型仅适用于自由扩散，无明显对流、受限空间与非理想梯度波形。
\item 数值回代核验残差：0.00e+0、0.00e+0、0.00e+0。
\end{enumerate}
\noindent\textbf{答案摘要}\par
答案1 = 1.23473473e+09 s/m²；答案2 = 0.476712909 signal fraction；答案3 = 6e-10 m²/s；答案4 = 0.0516449249 signal fraction\par
\noindent\textbf{模型前提：}
仅按题设理想化模型处理；实际化学性、跨辑核心模型独立性、版权和专家审校仍待核准。\par
\noindent\textcolor{red!65!black}{[跨辑去重未完成 · 草稿]} \par
\medskip\hrule\medskip
\Needspace{7\baselineskip}
\subsubsection*{UC-ADV10-A102 · EPR两个不等价核自旋的六重超精细分裂与谱线重心}
\noindent\textbf{题目}\par
一个假想有机自由基的S=1/2电子自旋耦合两个互不等价核：一个14N核自旋I\_N=1、各向同性超精细常数A\_N=15.00G，另一个1H核I\_H=1/2、常数A\_H=5.000G。假设高场一阶近似、核跃迁不改变mI、无可分辨二阶项及自旋驰豫差异，未分裂电子共振磁场B0=340.0mT，1G=0.1000mT。共振位置B=B0−mN A\_N−mH A\_H。求总谱线数、所有六条谱线位置及相对积分强度、线谱重心与总跨度，验证核态等简并时积分强度为1:1:1:1:1:1。\par
\medskip\noindent\textbf{分步解答}
\begin{enumerate}
\item 14N的核磁量子数mN=−1,0,+1有三个取值，1H的mH=−1/2,+1/2有两个取值。
\item 不等价核的高场一级超精细谱线总数(2I\_N+1)(2I\_H+1)=3×2=6。
\item 将A\_N=1.500mT、A\_H=0.5000mT代入B=B0−mN A\_N−mH A\_H，并排序得到338.25、338.75、339.75、340.25、341.25、341.75mT。
\item 假设每个核自旋投影等概率且所有六条互不重叠，每条积分信号等于总积分的1/6。
\item 对六条谱线等权求重心，Bcentroid=340mT。
\item 最大与最小场的跨度ΔB=3.5mT；这是核超精细展宽而非额外电子自旋态数。
\item 若两核A值使谱线完全重合，会出现组合重叠强度，需要按mN和mH逐线相加，而不能仍认作六条可分辨峰。
\item 独立代回原始配分、强度或热力学公式所得主要残差：0.000e+0、0.000e+0、0.000e+0；若超出题设近似则需重新建模。
\end{enumerate}
\noindent\textbf{答案摘要}\par
答案1 = 6 lines；答案2 = 338.25 mT；答案3 = 338.75 mT；答案4 = 339.75 mT；答案5 = 340.25 mT；答案6 = 341.25 mT；答案7 = 341.75 mT；答案8 = 340 mT；答案9 = 3.5 mT\par
\noindent\textbf{模型前提：}
标准态、理想活度和模型参数仅按本题设定成立，未经专家科学性及跨辑原创性审定。\par
\noindent\textcolor{red!65!black}{[跨辑去重未完成 · 草稿]} \par
\medskip\hrule\medskip
\Needspace{7\baselineskip}
\subsubsection*{UC-ADV10-A103 · 色谱峰方差的柱外展宽扣除与分离度真实测量}
\noindent\textbf{题目}\par
某高效液相色谱系统的柱长1.200m，峰1保留时间tR1=8.000min、检测到的高斯峰标准差σobs,1=0.2500min。独立使用零保留标记得到柱外装置标准差σext=0.1500min，假设柱内与柱外的高斯方差独立可加、无峰拖尾。第二峰tR2=8.800min、σobs,2=0.3000min，共用相同柱外展宽。求峰1柱内真实σ、表观及扣除柱外后的理论塔板数N=(tR/σ)\textasciicircum{}2、校正后塔板高度H=L/N、基于各峰σ的检测端表观分离度及剔除柱外展宽后的柱本征分离度，说明后者不是实际检测器已经实现的分离度。\par
\medskip\noindent\textbf{分步解答}
\begin{enumerate}
\item 高斯独立展宽必须以方差相加，σcol,1=√(σobs,1²−σext²)=0.2min，而不能直接从标准差做差。
\item 不扣除柱外展宽的表观塔板数Nobs=(8/0.25)²=1024。
\item 柱本征塔板数Ncol=(8/σcol,1)²=1600。
\item 由柱长1.200m得到Hcol=L/Ncol=0.00075m。
\item 峰2柱本征标准差σcol,2=√(0.3000²−0.1500²)=0.259807621min。
\item 以基线宽约4σ的公式Rs=Δt/[2(σ1+σ2)]，检测端Rs,obs=0.727272727，去柱外展宽的柱内Rs,col=0.869929035。
\item 柱内Rs描述理想去除外部散布后的本征性能；实际检测系统的峰仍具有柱外展宽，必须使用Rs,obs评估实测分离。
\item 独立代回原始配分、强度或热力学公式所得主要残差：0.000e+0、0.000e+0、0.000e+0；若超出题设近似则需重新建模。
\end{enumerate}
\noindent\textbf{答案摘要}\par
答案1 = 0.2 min；答案2 = 1024 plates；答案3 = 1600 plates；答案4 = 0.00075 m；答案5 = 0.259807621 min；答案6 = 0.727272727 ratio；答案7 = 0.869929035 ratio\par
\noindent\textbf{模型前提：}
标准态、理想活度和模型参数仅按本题设定成立，未经专家科学性及跨辑原创性审定。\par
\noindent\textcolor{red!65!black}{[跨辑去重未完成 · 草稿]} \par
\medskip\hrule\medskip
\Needspace{7\baselineskip}
\subsubsection*{UC-ADV10-A104 · 双指数荧光寿命的频域相量与强度加权均值}
\noindent\textbf{题目}\par
某假想时间分辨荧光样品受到理想瞬时激发，已去除仪器响应及背景。荧光衰减为I(t)=A₁exp(−t/τ₁)+A₂exp(−t/τ₂)，其中A₁=0.6500、A₂=0.3500，τ₁=1.500 ns、τ₂=5.000 ns，t≥0。定义积分光子份额f₂=∫I₂dt/∫Idt、平均光子到达时间〈t〉=∫tI dt/∫I dt，以及100.0 MHz调制角频率ω=2πf下的相量G=∫I cos(ωt)dt/∫I dt、S=∫I sin(ωt)dt/∫I dt。求f₂、〈t〉、G、S与按单指数关系误推的τ\_app=S/(ωG)，解释A₂为何不等于f₂。\par
\medskip\noindent\textbf{分步解答}
\begin{enumerate}
\item 将100.0 MHz写成0.1000 ns⁻¹，故ω=0.6283185307 rad/ns；积分变量全用ns，避免把MHz与ns直接相乘。
\item 单指数积分∫₀∞exp(−t/τ)dt=τ，因而总积分D=A₁τ₁+A₂τ₂=0.6500×1.500+0.3500×5.000=2.725 ns。
\item 第二组分积分光子份额f₂=A₂τ₂/D=0.6422018349；初始振幅A₂=0.3500并非积分份额，因为长寿命组分每单位初始振幅贡献更多光子。
\item 利用∫₀∞t exp(−t/τ)dt=τ²，得到〈t〉=(A₁τ₁²+A₂τ₂²)/D=3.747706422 ns。
\item 余弦积分为τ/[1+(ωτ)²]，所以G=Σ[A\_iτ\_i/(1+(ωτ\_i)²)]/D=0.2485675603。
\item 正弦积分为ωτ²/[1+(ωτ)²]，所以S=Σ[A\_iωτ\_i²/(1+(ωτ\_i)²)]/D=0.3641982929。
\item 若强行套用单指数相量式，τ\_app=S/(ωG)=2.331919681 ns，与真实强度加权平均〈t〉不同，不能将混合衰减当作单指数。
\item 积分归一化检查f₁+f₂=1；G、S均位于[0,1]且G²+S²<1。若存在散射、响应卷积或能量转移，必须另行修正，不能照搬本题相量。
\end{enumerate}
\noindent\textbf{答案摘要}\par
答案1 = 0.642201835 fraction；答案2 = 3.74770642 ns；答案3 = 0.24856756 dimensionless；答案4 = 0.364198293 dimensionless；答案5 = 2.33191968 ns\par
\noindent\textbf{模型前提：}
双指数独立辐射衰减、无仪器响应和背景、固定预指数系数；频率采用每ns单位，尚未专家审定。\par
\noindent\textcolor{red!65!black}{[跨辑去重未完成 · 草稿]} \par
\medskip\hrule\medskip
\subsection*{有机化学}
\addcontentsline{toc}{subsection}{第10辑 / 有机化学}
\Needspace{7\baselineskip}
\subsubsection*{UC-ADV10-O98 · 碳同位素交叉实验逆推分子间随机重组途径比例}
\noindent\textbf{题目}\par
为区分某假想有机重排中的分子内路径与分子间路径，将等摩尔的两类双片段底物混合：一类两个片段均带13C标记(L/L)，另一类两个片段均不标记(U/U)。分子内过程保留原配对，因此L/L和U/U各占1/2；分子间过程随机重组并假设无同位素效应，L/L、混合L/U、U/U的理论比例为1/4、1/2、1/4。实验生成总产物0.01000mol，测得L/U混合标记产物占30.00\%。求分子间途径比例、三类产物比例、混合标记产物mol和分子内途径贡献的产物mol；解释随机配对假设的限制。\par
\medskip\noindent\textbf{分步解答}
\begin{enumerate}
\item 设α为分子间随机重组贡献的总产物比例，其余1−α来自分子内途径。
\item 只有分子间途径形成混合L/U，故P(L/U)=α·2p(1−p)=α/2=0.3000，解α=0.6。
\item L/L概率=(1−α)/2+αp²=0.35，U/U同理=0.35。
\item 混合L/U比例=0.3，三项之和=1。
\item 混合产物摩尔数=0.01000·0.3000=0.003mol。
\item 分子内途径生成产物=(1−α)·0.01000=0.004mol。
\item 仅当分子间片段随机重组、产物检测灵敏度已校正且无笼内相关重组时，交叉比例才能反推宏观途径贡献。
\item 数值回代核验残差：0.00e+0、0.00e+0、0.00e+0。
\end{enumerate}
\noindent\textbf{答案摘要}\par
答案1 = 0.6 fraction；答案2 = 0.35 fraction；答案3 = 0.3 fraction；答案4 = 0.35 fraction；答案5 = 0.003 mol；答案6 = 0.004 mol\par
\noindent\textbf{模型前提：}
仅按题设理想化模型处理；实际化学性、跨辑核心模型独立性、版权和专家审校仍待核准。\par
\noindent\textcolor{red!65!black}{[跨辑去重未完成 · 草稿]} \par
\medskip\hrule\medskip
\Needspace{7\baselineskip}
\subsubsection*{UC-ADV10-O99 · 氢氘动力学同位素效应的Eyring焓熵拆分与温度反演}
\noindent\textbf{题目}\par
假想有机C–H键断裂为速率决定步骤的反应，氘代同系物在相同溶剂下形成同一类型的产物，给定实测一级动力学同位素效应KIE=kH/kD：298.15K为6.000、340.0K为4.000。假设在这一温度区间两者过渡态活化焓差ΔΔH‡=ΔH‡H−ΔH‡D与活化熵差ΔΔS‡=ΔS‡H−ΔS‡D恒定，且Eyring通用kBT/h因子在速率比中相消。由ln(kH/kD)=−ΔΔH‡/(RT)+ΔΔS‡/R反推出ΔΔH‡、ΔΔS‡、298.15K的ΔΔG‡，预测320K时KIE，并解释为何较低温的KIE更大并不能直接证明量子隧穿。\par
\medskip\noindent\textbf{分步解答}
\begin{enumerate}
\item 同温度速率比消去Eyring共同kBT/h项，得到lnKIE对1/T的直线，斜率m=−ΔΔH‡/R。
\item 两温度点直线斜率m=[ln4−ln6]/[1/340−1/298.15]=982.136284K。
\item 因此ΔΔH‡=−R m=-8.16593542kJ/mol；符号负表示H通道相对D有更低活化焓。
\item 由298.15K的ln6=−ΔΔH‡/(RT1)+ΔΔS‡/R，求得ΔΔS‡=-12.4911645J/(mol K)。
\item 298.15K时ΔΔG‡=ΔΔH‡−T1ΔΔS‡=-4.44169473kJ/mol。
\item 在320K按恒焓熵关系预测KIE=4.79145498，介于4和6之间。
\item 普通零点能差即可造成温度依赖KIE；隧穿需要更多温度点、基团结构和动力学证据，不能由两点比值单独证成。
\item 由题设能量、物料或统计分布公式重代所得归一化残差：0.000e+0、-6.661e-16、0.000e+0。
\end{enumerate}
\noindent\textbf{答案摘要}\par
答案1 = -8.16593542 kJ/mol；答案2 = -12.4911645 J/(mol K)；答案3 = -4.44169473 kJ/mol；答案4 = 4.79145498 ratio；答案5 = 982.136284 K\par
\noindent\textbf{模型前提：}
理想状态、独立过程、纯一阶模型按题设成立；实验事实和跨辑模型独立性待专家确认。\par
\noindent\textcolor{red!65!black}{[跨辑去重未完成 · 草稿]} \par
\medskip\hrule\medskip
\Needspace{7\baselineskip}
\subsubsection*{UC-ADV10-O100 · 有机光脱保护薄膜的Beer–Lambert吸收深度分布与量子转化}
\noindent\textbf{题目}\par
一层均匀有机光脱保护聚合物膜厚L=100.0μm，采用365nm单色光照射、入射光子通量F0=1.000×10\textasciicircum{}20 photons/(m²s)，介质单程吸收系数α=2.000×10\textasciicircum{}4m\textasciicircum{}-1。假设没有反射、散射、二次吸收与产物遮光，局部光子通量F(z)=F0exp(−αz)，每吸收一个光子生成目标脱保护产物的量子产率φ=0.1500，产物均留在膜中。求全膜光子吸收分数、每平方米每秒的目标产物mol生成速率、全部被吸收光子的位置平均深度，以及前半层(0至L/2)占所有光子吸收的比例；说明为何不能把薄膜吸收当作沿厚度均匀分布。\par
\medskip\noindent\textbf{分步解答}
\begin{enumerate}
\item 光学厚度αL=2，全膜的光子吸收率A=1−exp(−αL)=0.864664717。
\item 单位面积每秒吸收光子数Fabs=F0A=86466471700000000000photons/(m²s)。
\item 目标产物摩尔通量Jproduct=φFabs/NA=0.0000215371431mol/(m²s)。
\item 沿膜厚的光子吸收分布密度正比于αexp(−αz)，不能用位置均匀分布计算平均深度。
\item 条件于光子确已被全膜吸收，平均吸收深度〈z〉=[1−(1+αL)exp(−αL)]/[α(1−exp(−αL))]=34.3482357μm。
\item 前半膜吸收份额fhalf=[1−exp(−αL/2)]/[1−exp(−αL)]=0.731058579，高于均匀膜的0.5。
\item 极限αL趋零时平均吸收深度趋L/2、前半膜份额趋1/2；极限αL很大时吸收趋向入射表面。
\item 由题设能量、物料或统计分布公式重代所得归一化残差：0.000e+0、1.110e-16、0.000e+0。
\end{enumerate}
\noindent\textbf{答案摘要}\par
答案1 = 0.864664717 fraction；答案2 = 2.15371431e-05 mol/(m²s)；答案3 = 34.3482357 μm；答案4 = 0.731058579 fraction；答案5 = 2 dimensionless\par
\noindent\textbf{模型前提：}
理想状态、独立过程、纯一阶模型按题设成立；实验事实和跨辑模型独立性待专家确认。\par
\noindent\textcolor{red!65!black}{[跨辑去重未完成 · 草稿]} \par
\medskip\hrule\medskip
\Needspace{7\baselineskip}
\subsubsection*{UC-ADV10-O101 · 动态动力学拆分中消旋互变与双通道产物对映选择性}
\noindent\textbf{题目}\par
假想手性底物的R、S对映体在恒温均相体系中以相同一级速率常数k\_r=0.0400 min⁻¹双向消旋互变R⇌S；同时R以k\_R=0.1200 min⁻¹不可逆生成P\_R，S以k\_S=0.0200 min⁻¹不可逆生成P\_S。初始总底物归一化为1.0000 mol，其中R₀=S₀=0.5000 mol，P\_R,0=P\_S,0=0。反应体积恒定，所有通道一级、无副反应，反应20.00 min。求剩余R、S摩尔数、两产物摩尔数、产物对映体过量ee\_P=(P\_R−P\_S)/(P\_R+P\_S)及转化率，并核对质量守恒与消旋极限的意义。\par
\medskip\noindent\textbf{分步解答}
\begin{enumerate}
\item 设底物列向量v=(R,S)\textasciicircum{}T，则dv/dt=Mv，M=[[-(k\_R+k\_r),k\_r],[k\_r,−(k\_S+k\_r)]]，初始v₀=(0.5,0.5)\textasciicircum{}T。产物满足dP\_R/dt=k\_R R、dP\_S/dt=k\_S S。
\item 令m=(k\_R+k\_S+2k\_r)/2=0.1100 min⁻¹，d=√\{[(k\_R−k\_S)/2]²+k\_r²\}=0.0640312424 min⁻¹，矩阵指数为e\textasciicircum{}(Mt)=e\textasciicircum{}(−mt)[cosh(dt)I+(sinh(dt)/d)(M+mI)]。
\item 由(M+mI)v₀=(−0.0050,0.0450)\textasciicircum{}T min⁻¹，故R(t)=e\textasciicircum{}(−mt)[0.5cosh(dt)−0.005sinh(dt)/d]，S(t)=e\textasciicircum{}(−mt)[0.5cosh(dt)+0.045sinh(dt)/d]。
\item t=20.00 min时，R=0.09302187103 mol、S=0.23669390978 mol；因R被优先消耗，剩余底物富集S。
\item 设J\_R=∫₀ᵗR(u)du、J\_S=∫₀ᵗS(u)du。积分微分方程给M(J\_R,J\_S)\textasciicircum{}T=v(t)−v₀，因det M=(k\_R+k\_r)(k\_S+k\_r)−k\_r²=0.0080 min⁻²>0，可唯一反解。
\item 由P\_R=k\_R J\_R、P\_S=k\_S J\_S，得到P\_R=0.52426397020 mol、P\_S=0.14602024899 mol。不能简单以瞬时R/S比例乘以速率得到累计产物。
\item 产物ee\_P=(P\_R−P\_S)/(P\_R+P\_S)=0.5643034856，转化率X=P\_R+P\_S=0.6702842192；R+S+P\_R+P\_S=1.0000000000 mol。
\item 若k\_r→0则两底物独立指数衰减；若k\_r远大于两个反应速率，底物近似快速消旋，产物分流趋向k\_R:k\_S。本题无外消旋副反应与产物再互变。
\end{enumerate}
\noindent\textbf{答案摘要}\par
答案1 = 0.093021871 mol；答案2 = 0.23669391 mol；答案3 = 0.52426397 mol；答案4 = 0.146020249 mol；答案5 = 0.564303486 fraction；答案6 = 0.670284219 fraction\par
\noindent\textbf{模型前提：}
封闭体系，双向消旋速率对称且产物不可逆、不互变；速率常数恒定；模型新颖性与实际机理待专家审校。\par
\noindent\textcolor{red!65!black}{[跨辑去重未完成 · 草稿]} \par
\medskip\hrule\medskip
\subsection*{物理化学}
\addcontentsline{toc}{subsection}{第10辑 / 物理化学}
\Needspace{7\baselineskip}
\subsubsection*{UC-ADV10-P98 · Bosanquet分子扩散与Knudsen孔壁阻力耦合的膜通量}
\noindent\textbf{题目}\par
一层假想多孔氮气扩散膜在320.0K工作，孔半径r=100.0nm，氮气摩尔质量M=0.02800kg/mol、自由气体扩散系数Dm=1.800×10\textasciicircum{}-5m²/s。孔壁碰撞的Knudsen扩散系数DK=(2r/3)√(8RT/(πM))，孔内总扩散遵循Bosanquet阻力式1/Dp=1/Dm+1/DK。膜孔隙率ε=0.3500、曲折度τ=2.500、厚度L=0.8000mm，膜两面浓度分别c1=3.000mol/m³、c2=0.5000mol/m³。求DK、Dp、膜有效Deff=εDp/τ、稳态通量及忽略Knudsen阻力时通量高估倍数。\par
\medskip\noindent\textbf{分步解答}
\begin{enumerate}
\item 半径和厚度转为SI：r=1×10\textasciicircum{}-7m，L=8×10\textasciicircum{}-4m。
\item Knudsen系数DK=(2r/3)√(8RT/(πM))=0.00003279380909m²/s。
\item 分子与孔壁散射阻力串联，Dp=[1/Dm+1/DK]\textasciicircum{}(-1)=0.00001162126988m²/s。
\item 实际多孔膜Deff=εDp/τ=0.000001626977783m²/s。
\item Fick稳态通量J=Deff(c1−c2)/L=0.005084305572mol/(m²s)。
\item 遗漏孔壁阻力会错误获得Jwrong=0.007875mol/(m²s)，高估倍率=1.548884088。
\item r→∞时DK→∞，Dp→Dm；r→0时孔壁Knudsen过程控制，且该模型没有计入表面吸附扩散。
\item 数值回代核验残差：0.00e+0、0.00e+0、0.00e+0。
\end{enumerate}
\noindent\textbf{答案摘要}\par
答案1 = 3.27938091e-05 m²/s；答案2 = 1.16212699e-05 m²/s；答案3 = 1.62697778e-06 m²/s；答案4 = 0.00508430557 mol/(m²s)；答案5 = 1.54888409 ratio\par
\noindent\textbf{模型前提：}
仅按题设理想化模型处理；实际化学性、跨辑核心模型独立性、版权和专家审校仍待核准。\par
\noindent\textcolor{red!65!black}{[跨辑去重未完成 · 草稿]} \par
\medskip\hrule\medskip
\Needspace{7\baselineskip}
\subsubsection*{UC-ADV10-P99 · 单级热电离Saha平衡的正根选择与等温压强}
\noindent\textbf{题目}\par
某假想稀薄单原子气体在T=8000K达到局部热平衡，只发生第一电离A⇌A++e−，忽略二次电离、电子简并与粒子间强相互作用。气体原子核总数密度N=nA+nA+=1.000×10\textasciicircum{}21m\textasciicircum{}-3，已把电子及原子内部配分因素合并的Saha系数S(T)=4.000×10\textasciicircum{}20m\textasciicircum{}-3，给定平衡式n\_e n\_A+/n\_A=S(T)，并满足电中性n\_e=n\_A+。求电离分数x=nA+/N、电子与中性原子数密度、理想总压强p=(nA+nA++n\_e)kBT，以及若每个电离原子储存12.00eV电离能时每立方米的电离内能；严格选择0≤x≤1的正根。\par
\medskip\noindent\textbf{分步解答}
\begin{enumerate}
\item 定义电离分数x=nA+/N，电中性给n\_e=Nx，且nA=N(1−x)。
\item Saha平衡写为x²/(1−x)=S/N=0.4，整理为x²+Kx−K=0。
\item 物理解取x=[−K+√(K²+4K)]/2=0.463324958，另一负根不满足数密度非负性。
\item 电子数密度ne=Nx=463324958000000000000m\textasciicircum{}-3，中性原子数密度nA=N(1−x)=536675042000000000000m\textasciicircum{}-3。
\item 总粒子数密度N+ne=N(1+x)，理想压强p=N(1+x)kBT=161.627051Pa。
\item 若每次第一电离固定耗能12.00eV，电离内能密度u\_ion=ni×12eV=890.794106J/m³。
\item S(T)本身随温度与粒子配分函数变化，本题只使用已给定T的系数；非理想高密度等离子体还需校正电离能降低等效应。
\item 由题设能量、物料或统计分布公式重代所得归一化残差：2.220e-16、0.000e+0、0.000e+0。
\end{enumerate}
\noindent\textbf{答案摘要}\par
答案1 = 0.463324958 fraction；答案2 = 4.63324958e+20 m\textasciicircum{}-3；答案3 = 5.36675042e+20 m\textasciicircum{}-3；答案4 = 161.627051 Pa；答案5 = 890.794106 J/m³\par
\noindent\textbf{模型前提：}
理想状态、独立过程、纯一阶模型按题设成立；实验事实和跨辑模型独立性待专家确认。\par
\noindent\textcolor{red!65!black}{[跨辑去重未完成 · 草稿]} \par
\medskip\hrule\medskip
\Needspace{7\baselineskip}
\subsubsection*{UC-ADV10-P100 · 热膨胀和等温压缩率决定Cp−Cv及PVT偏导的Maxwell闭合}
\noindent\textbf{题目}\par
某假想单相液体在300.0K和当前参考压强处有摩尔体积V=1.800×10\textasciicircum{}-5m³/mol、定压摩尔热容Cp=75.00J/(mol K)、体膨胀系数α=(1/V)(∂V/∂T)P=2.000×10\textasciicircum{}-4K\textasciicircum{}-1和等温压缩率κT=−(1/V)(∂V/∂P)T=5.000×10\textasciicircum{}-10Pa\textasciicircum{}-1，假设小范围内这些参数不变。由热力学恒等式Cp−Cv=TVα²/κT与(∂P/∂T)V=α/κT求定容热容Cv、定容升温压强斜率、定压体积温度导数和等温体积压强导数；另用常Cp近似算300至310K定压过程熵增，分别检查量纲和正稳定性。\par
\medskip\noindent\textbf{分步解答}
\begin{enumerate}
\item 根据热力学恒等式Cp−Cv=T V α²/κT=0.432J/(mol K)，它因κT>0而非负。
\item Cv=Cp−(Cp−Cv)=74.568J/(mol K)，仍应严格大于0。
\item 定容压强温度导数(∂P/∂T)V=α/κT=400000Pa/K。
\item 定压体积温度导数(∂V/∂T)P=αV=3.6e-9m³/(mol K)。
\item 等温体积压强导数(∂V/∂P)T=−κTV=-9e-15m³/(mol Pa)，符号必须为负。
\item 在Cp近似恒定下300至310K定压过程的摩尔熵增ΔS=Cp ln(310/300)=2.45923671J/(mol K)。
\item 统一以SI处理热容差，检查Cp>Cv>0及κT>0；若跨越相变或参数强烈变化，应积分真实状态函数而非外推这些常数。
\item 由题设能量、物料或统计分布公式重代所得归一化残差：2.109e-15、0.000e+0、0.000e+0。
\end{enumerate}
\noindent\textbf{答案摘要}\par
答案1 = 74.568 J/(mol K)；答案2 = 0.432 J/(mol K)；答案3 = 400000 Pa/K；答案4 = 3.6e-09 m³/(mol K)；答案5 = -9e-15 m³/(mol Pa)；答案6 = 2.45923671 J/(mol K)\par
\noindent\textbf{模型前提：}
理想状态、独立过程、纯一阶模型按题设成立；实验事实和跨辑模型独立性待专家确认。\par
\noindent\textcolor{red!65!black}{[跨辑去重未完成 · 草稿]} \par
\medskip\hrule\medskip
\Needspace{7\baselineskip}
\subsubsection*{UC-ADV10-P101 · 弱酸固定电荷水凝胶的Donnan分配与电荷调节耦合}
\noindent\textbf{题目}\par
在298.15 K，一理想带弱酸基团的水凝胶与无限大NaCl水溶液接触。外液NaCl浓度c\_s=0.1000 mol/L，外液pH=7.000由极微量缓冲体系固定；凝胶固定酸位点总浓度C\_T=0.1500 mol/L凝胶水相，pKa=6.500，解离为不可移动A⁻。Na⁺、Cl⁻、H⁺可透膜，固定酸基团不可透膜；忽略H⁺/OH⁻对电中性的微小贡献和活度系数。令r=[Na⁺]\_in/c\_s，理想Donnan分配满足[Cl⁻]\_in=c\_s/r、[H⁺]\_in=r[H⁺]\_out。求r、内部Na⁺与Cl⁻浓度、固定酸位点解离度α、内部pH、内外电势差ψ\_in−ψ\_out及仅由可移动NaCl离子产生的渗透压差ΔΠ\_ion。\par
\medskip\noindent\textbf{分步解答}
\begin{enumerate}
\item 由同价阳离子电化学势相等，[Na⁺]\_in=c\_s r且[H⁺]\_in=r[H⁺]\_out；阴离子反向分配[Cl⁻]\_in=c\_s/r，故pH\_in=pH\_out−log₁₀r。
\item 固定酸位点服从Henderson–Hasselbalch关系α=1/[1+10\textasciicircum{}(pKa−pH\_in)]=1/(1+a r)，其中a=10\textasciicircum{}(6.500−7.000)=0.3162277660。
\item 忽略微量H⁺/OH⁻电荷后，凝胶内部电中性要求c\_s(r−1/r)=C\_T α。乘以r(1+a r)得到三次方程a c\_s r³+c\_s r²−(a c\_s+C\_T)r−c\_s=0。
\item 左侧净正离子c\_s(r−1/r)在r≥1严格递增，右侧固定负电荷C\_T/(1+a r)严格递减，故物理解唯一且r>1。数值解r=1.6131825603。
\item 内部Na⁺浓度0.1613182560 mol/L、Cl⁻浓度0.06198926424 mol/L，差0.09932899179 mol/L，恰等于C\_T α。
\item 解离度α=0.6621932786，内部pH=6.792316482；凝胶内部酸性高于外液，因此固定酸位点部分被抑制解离。
\item 阳离子电化学势关系给ψ\_in−ψ\_out=−(RT/F)ln r=−0.01228642188 V；内部电势较负与阳离子富集相符。
\item 仅计移动NaCl离子，ΔΠ\_ion=RT×1000×(c\_s r+c\_s/r−2c\_s)=57778.34122 Pa，mol/L需乘1000变为mol/m³。固定聚合物网络的混合与弹性渗透压未计入，故此值不是完整凝胶溶胀压。
\item 检查电中性残差|c\_Na−c\_Cl−C\_Tα|<10⁻¹² mol/L、Donnan乘积c\_Na c\_Cl=c\_s²及r>1；若离子活度非理想或缓冲离子浓度不可忽略，需要联立更多物种。
\end{enumerate}
\noindent\textbf{答案摘要}\par
答案1 = 1.61318256 ratio；答案2 = 0.161318256 mol/L；答案3 = 0.0619892642 mol/L；答案4 = 0.662193279 fraction；答案5 = 6.79231648 pH；答案6 = -0.0122864219 V；答案7 = 57778.3412 Pa\par
\noindent\textbf{模型前提：}
固定酸基团不贡献移动粒子渗透压；理想单价盐，外部缓冲离子忽略，忽略H⁺/OH⁻电中性贡献，实际凝胶溶胀仍需其他能量项。\par
\noindent\textcolor{red!65!black}{[跨辑去重未完成 · 草稿]} \par
\medskip\hrule\medskip
\end{document}
